Con el objetivo de cubrir una región de interés agrícola, el análisis se centró en los estados federados de Baviera, Brandemburgo y Sajonia, los cuales, de manera conjunta, albergan 5.17 M ha de superficie destinada para la agricultura \citep{Flores_2025}. Esta área representa aproximadamente el $54.42\%$ de las 9.5 M ha de superficie empleada anualmente para los principales cultivos que se producen en Alemania \citep{Duden_2024}.\\


Para cada uno de los estados mencionados se replicó, de manera independiente, el criterio empírico de permanencia antes descrito para generar los correspondientes calendarios agrícolas regionalizados para la siembra de maíz en Alemania. Para ello, se consideraron exclusivamente los valores registrados diariamente dentro de las celdas contenidas en sus respectivas fronteras geográficas, ver Figura \ref{Stations_per_state}.\\


Para integrar los efectos del cambio climático se consideró un periodo de 30 años de duración, un tiempo estándar en el área \citep{Hennemuth_2013}. Por tanto se contemplaron los datos históricos desde 1996 hasta 2025. Para cada variable considerada y para cada día del año, se calculó el promedio de los valores reportados en las celdas  contenidas en el estado correspondiente.\\

De esta manera, se obtienen mediciones históricas representativas de las condiciones medias locales, ver Figura \ref{Distance}, las cuales son utilizadas para determinar las probabilidades de transición del modelo a través de la Ecuación \eqref{lambdas}.\\ 


\begin{figure}[!h]
	\centering
	\includegraphics[width=0.5\textwidth]{Mapa_final.png}
	\caption{Distribución de las estaciones climáticas respecto a los estados federados de Alemania. Los polígonos delimitados por un contorno en color negro corresponden a los estados federados reportados por Bundesamt für Kartographie und Geodäsie \cite{States_2025}. En color verde se muestran, de manera descendente, Brandemburgo, Sajonia y Baviera, respectivamente.}
	\label{Stations_per_state}
\end{figure}


Con lo anterior, para cada estado, se obtuvieron un par de calendarios agrícolas: uno empleando la distancia Euclidiana y uno empleando la distancia Manhattan. Lo que permitió la identificación de periodos que favorezcan la siembra, adaptados a las condiciones locales mediante los criterios de decisión previamente establecidos, los cuales se muestran en la Figura \ref{fig:sowing_suitability}


% ============================================================
% Brandenburgo
% ============================================================

\begin{figure}[htbp]
	\centering
	\begin{subfigure}[b]{0.43\textwidth}
		\centering
		\includegraphics[width=\textwidth]{[Brandenburg][Euclidean]Transition_graph.png}
		\caption{}
		\label{fig:brandenburg_euclidean_graph}
	\end{subfigure}
	\hspace{0.03\textwidth}
	\begin{subfigure}[b]{0.43\textwidth}
		\centering
		\includegraphics[width=\textwidth]{[Brandenburg][Manhattan]Transition_graph.png}
		\caption{}
		\label{fig:brandenburg_manhattan_graph}
	\end{subfigure}
	
	
	\caption{
		Aptitud diaria para la siembra en Brandenburgo, expresada como la probabilidad de transición $\lambda_k$, $k\in\{1,2,\ldots,365\}$, calculada mediante las distancias Euclidiana y Manhattan. La figura \ref{fig:brandenburg_euclidean_graph} muestra los valores de $\lambda_k$ calculados mediante la distancia Euclidiana, mientras que la figura \ref{fig:brandenburg_manhattan_graph} muestra los valores correspondientes calculados mediante la distancia Manhattan. La progresión del año se expresa en sentido contrario a las manecillas del reloj; en el exterior de la gráfica anular se indican los meses, mientras que en su interior se indican los días del año.
	}
	\label{fig:sowing_suitability_brandenburg}
	
	
\end{figure}

% ============================================================
% Sajonia
% ============================================================

\begin{figure}[htbp]
	\centering
	\begin{subfigure}[b]{0.43\textwidth}
		\centering
		\includegraphics[width=\textwidth]{[Saxony][Euclidean]Transition_graph.png}
		\caption{}
		\label{fig:saxony_euclidean_graph}
	\end{subfigure}
	\hspace{0.03\textwidth}
	\begin{subfigure}[b]{0.43\textwidth}
		\centering
		\includegraphics[width=\textwidth]{[Saxony][Manhattan]Transition_graph.png}
		\caption{}
		\label{fig:saxony_manhattan_graph}
	\end{subfigure}
	
	
	\caption{
		Aptitud diaria para la siembra en Sajonia, expresada como la probabilidad de transición $\lambda_k$, $k\in\{1,2,\ldots,365\}$, calculada mediante las distancias Euclidiana y Manhattan. La figura \ref{fig:saxony_euclidean_graph} muestra los valores de $\lambda_k$ calculados mediante la distancia Euclidiana, mientras que la figura \ref{fig:saxony_manhattan_graph} muestra los valores correspondientes calculados mediante la distancia Manhattan. La progresión del año se expresa en sentido contrario a las manecillas del reloj; en el exterior de la gráfica anular se indican los meses, mientras que en su interior se indican los días del año.
	}
	\label{fig:sowing_suitability_saxony}
	
	
\end{figure}

% ============================================================
% Baviera
% ============================================================

\begin{figure}[htbp]
	\centering
	\begin{subfigure}[b]{0.43\textwidth}
		\centering
		\includegraphics[width=\textwidth]{[Bavaria][Euclidean]Transition_graph.png}
		\caption{}
		\label{fig:bavaria_euclidean_graph}
	\end{subfigure}
	\hspace{0.03\textwidth}
	\begin{subfigure}[b]{0.43\textwidth}
		\centering
		\includegraphics[width=\textwidth]{[Bavaria][Manhattan]Transition_graph.png}
		\caption{}
		\label{fig:bavaria_manhattan_graph}
	\end{subfigure}
	
	
	\caption{
		Aptitud diaria para la siembra en Baviera, expresada como la probabilidad de transición $\lambda_k$, $k\in\{1,2,\ldots,365\}$, calculada mediante las distancias Euclidiana y Manhattan. La figura \ref{fig:bavaria_euclidean_graph} muestra los valores de $\lambda_k$ calculados mediante la distancia Euclidiana, mientras que la figura \ref{fig:bavaria_manhattan_graph} muestra los valores correspondientes calculados mediante la distancia Manhattan. La progresión del año se expresa en sentido contrario a las manecillas del reloj; en el exterior de la gráfica anular se indican los meses, mientras que en su interior se indican los días del año.
	}
	\label{fig:sowing_suitability_bavaria}

	
\end{figure}

\iffalse
\begin{figure}[htbp]
	\centering
	\begin{tabular}{@{\hspace{-2cm}}c@{\hspace{-0.3cm}}c@{\hspace{-0.3cm}}c@{\hspace{0cm}}c}
		
		% ====================================================
		% Primera fila
		% ====================================================
		
		\begin{subfigure}[b]{0.38\textwidth}
			\centering
			\includegraphics[width=\textwidth]{[Brandenburg][Euclidean]Transition_graph.png}
			\caption{}
			\label{fig:brandenburg_euclidean_graph}
		\end{subfigure}
		&
		\begin{subfigure}[b]{0.38\textwidth}
			\centering
			\includegraphics[width=\textwidth]{[Brandenburg][Manhattan]Transition_graph.png}
			\caption{}
			\label{fig:brandenburg_manhattan_graph}
		\end{subfigure}
		\\[2mm]
		
		% ====================================================
		% Segunda fila
		% ====================================================
		
		\begin{subfigure}[b]{0.4\textwidth}
			\centering
			\includegraphics[width=\textwidth]{[Saxony][Euclidean]Transition_graph.png}
			\caption{}
			\label{fig:saxony_euclidean_graph}
		\end{subfigure}
		&
		\begin{subfigure}[b]{0.4\textwidth}
			\centering
			\includegraphics[width=\textwidth]{[Saxony][Manhattan]Transition_graph.png}
			\caption{}
			\label{fig:saxony_manhattan_graph}
		\end{subfigure}
		&
		\\[2mm]
		\begin{subfigure}[b]{0.4\textwidth}
			\centering
			\includegraphics[width=\textwidth]{[Bavaria][Euclidean]Transition_graph.png}
			\caption{}
			\label{fig:bavaria_euclidean_graph}
		\end{subfigure}
		&
		\begin{subfigure}[b]{0.4\textwidth}
			\centering
			\includegraphics[width=\textwidth]{[Bavaria][Manhattan]Transition_graph.png}
			\caption{}
			\label{fig:bavaria_manhattan_graph}
		\end{subfigure}
		&
	\end{tabular}
	
	\caption{
		Aptitud diaria para la siembra en Brandenburgo, Sajonia y Baviera, expresada como la probabilidad de transición $\lambda_k$, $k\in\{1,2,\ldots,365\}$, calculada mediante las distancias Euclidiana y Manhattan. Las figuras \ref{fig:brandenburg_euclidean}, \ref{fig:saxony_euclidean} y \ref{fig:bavaria_euclidean} muestran los valores de $\lambda_k$ calculados mediante la distancia Euclidiana para Brandenburgo, Sajonia y Baviera, respectivamente. Similarmente, las figuras \ref{fig:brandenburg_manhattan}, \ref{fig:saxony_manhattan} y \ref{fig:bavaria_manhattan} muestran los valores correspondientes calculados mediante la distancia Manhattan. La progresión del año se expresa en sentido contrario a las manecillas del reloj, en el exterior de la gráfica anular se indican los meses, mientras que en su interior se indican los días del año.
	}
	\label{fig:sowing_suitability}
\end{figure}
\fi




% ============================================================
% Brandenburgo
% ============================================================





\begin{figure}[htbp]
	\centering
	
	\begin{subfigure}[b]{0.46\textwidth}
		\centering
		\hspace{-0.7cm}\includegraphics[width=\textwidth]{[Brandenburg][Circular][Euclidean]Sowing_windows.png}
		\caption{}
		\label{fig:brandenburg_euclidean_sowing_window}
	\end{subfigure}%
	\hspace{-0.3cm}%
	\begin{subfigure}[b]{0.46\textwidth}
		\centering
		\includegraphics[width=\textwidth]{[Brandenburg][Circular][Manhattan]Sowing_windows.png}
		\caption{}
		\label{fig:brandenburg_manhattan_sowing_window}
	\end{subfigure}
	\hspace{-0.3cm}
	\begin{subfigure}[b]{0.09\textwidth}
		\centering
		\raisebox{1cm}{%
			\includegraphics[width=0.6\textwidth]{Sowing suitability.png}
		}
	\end{subfigure}
	\caption{
		Calendario de los periodos de siembra identificados para Brandenburgo mediante las distancias Euclidiana y Manhattan. Las figuras \ref{fig:brandenburg_euclidean_sowing_window} y \ref{fig:brandenburg_manhattan_sowing_window} muestran los periodos resultantes para cada distancia, respectivamente.
	}
	\label{fig:brandenburg_sowing_windows}
	
\end{figure}







% ============================================================
% Sajonia
% ============================================================



\begin{figure}[htbp]
	\centering
	\begin{subfigure}[b]{0.46\textwidth}
		\centering
		\hspace{-1cm}\includegraphics[width=\textwidth]{[Saxony][Circular][Euclidean]Sowing_windows.png}
		\caption{}
		\label{fig:saxony_euclidean_sowing_window}
	\end{subfigure}%
	\hspace{-0.3cm}%
	\begin{subfigure}[b]{0.46\textwidth}
		\centering
		\includegraphics[width=\textwidth]{[Saxony][Circular][Manhattan]Sowing_windows.png}
		\caption{}
		\label{fig:saxony_manhattan_sowing_window}
	\end{subfigure}
	\hspace{-0.3cm}
	\begin{subfigure}[b]{0.09\textwidth}
		\centering
		\raisebox{1cm}{%
			\includegraphics[width=0.6\textwidth]{Sowing suitability.png}
		}
	\end{subfigure}
	\caption{
		Calendario de los periodos de siembra identificados para Sajonia mediante las distancias Euclidiana y Manhattan. Las figuras \ref{fig:saxony_euclidean_sowing_window} y \ref{fig:saxony_manhattan_sowing_window} muestran los periodos resultantes para cada distancia, respectivamente.
	}
	\label{fig:saxony_sowing_windows}
	
\end{figure}

% ============================================================
% Baviera
% ============================================================




\begin{figure}[htbp]
	\centering
	
	\begin{subfigure}[b]{0.46\textwidth}
		\centering
		\hspace{-0.7cm}\includegraphics[width=1\textwidth]{[Bavaria][Circular][Euclidean]Sowing_windows.png}
		\caption{}
		\label{fig:bavaria_euclidean_sowing_window}
	\end{subfigure}%
	\hspace{-0.3cm}%
	\begin{subfigure}[b]{0.46\textwidth}
		\centering
		\includegraphics[width=1\textwidth]{[Bavaria][Circular][Manhattan]Sowing_windows.png}
		\caption{}
		\label{fig:bavaria_manhattan_sowing_window}
	\end{subfigure}
	\hspace{-0.3cm}
	\begin{subfigure}[b]{0.09\textwidth}
		\centering
		\raisebox{1cm}{%
			\includegraphics[width=0.6\textwidth]{Sowing suitability.png}
		}
	\end{subfigure}
	\caption{
		Calendario de los periodos de siembra identificados para Baviera mediante las distancias Euclidiana y Manhattan. Las figuras \ref{fig:bavaria_euclidean_sowing_window} y \ref{fig:bavaria_manhattan_sowing_window} muestran los periodos resultantes para cada distancia, respectivamente.
	}
	\label{fig:bavaria_sowing_windows}
	
\end{figure}



\iffalse


%################################
%								#
%	     BRANDENBURG 			#
%								#
%################################



\begin{figure}[htbp]
	\centering
	
	
	% ============================================================
	% Primera fila: Raw
	% ============================================================
	\begin{subfigure}[b]{0.8\textwidth}
	\includegraphics[
	width=\textwidth
	]{[Brandenburg][Euclidean][Raw]average_temperature.png}
	\caption{}
	\label{fig:brandenburg_euclidean_raw_yearly_temperature}
	\end{subfigure}
	\vspace{2mm}
	
	% ============================================================
	% Segunda fila: CI
	% ============================================================
	\begin{subfigure}[b]{0.8\textwidth}
	\includegraphics[
	width=\textwidth
	]{[Brandenburg][Euclidean][CI]average_temperature.png}
	\caption{}
	\label{fig:brandenburg_euclidean_ci_yearly_temperature}
	\end{subfigure}
	\vspace{2mm}
	
	% ============================================================
	% Tercera fila: datos agregados
	% ============================================================
	
	\begin{subfigure}[b]{0.4\textwidth}
		\centering
		\includegraphics[
		width=\textwidth
		]{[Brandenburg][Euclidean][Raw]Aggregated_average_temperature.png}
		\caption{}
		\label{fig:brandenburg_euclidean_raw_aggregated_temperature}
	\end{subfigure}
	\hfill
	\begin{subfigure}[b]{0.4\textwidth}
		\centering
		\includegraphics[
		width=\textwidth
		]{[Brandenburg][Euclidean][CI]Aggregated_average_temperature.png}
		\caption{}
		\label{fig:brandenburg_euclidean_ci_aggregated_temperature}
	\end{subfigure}
	
	\caption{
		Temperatura media del aire para Brandenburgo obtenida mediante la distancia Euclidiana. Las dos primeras representaciones muestran los resultados para los datos sin agregar y sus respectivos intervalos de confianza, mientras que las dos representaciones inferiores muestran los resultados agregados para los mismos casos.
	}
	\label{fig:brandenburg_euclidean_average_temperature}

	
\end{figure}



\begin{figure}[htbp]
	\centering
	
	
	% ============================================================
	% Primera fila: Raw
	% ============================================================
	\begin{subfigure}[b]{0.8\textwidth}
		\includegraphics[
		width=\textwidth
		]{[Brandenburg][Manhattan][Raw]average_temperature.png}
		\caption{}
		\label{fig:brandenburg_Manhattan_raw_yearly_temperature}
	\end{subfigure}
	\vspace{2mm}
	
	% ============================================================
	% Segunda fila: CI
	% ============================================================
	\begin{subfigure}[b]{0.8\textwidth}
		\includegraphics[
		width=\textwidth
		]{[Brandenburg][Manhattan][CI]average_temperature.png}
		\caption{}
		\label{fig:brandenburg_Manhattan_ci_yearly_temperature}
	\end{subfigure}
	\vspace{2mm}
	
	% ============================================================
	% Tercera fila: datos agregados
	% ============================================================
	
	\begin{subfigure}[b]{0.4\textwidth}
		\centering
		\includegraphics[
		width=\textwidth
		]{[Brandenburg][Manhattan][Raw]Aggregated_average_temperature.png}
		\caption{}
		\label{fig:brandenburg_Manhattan_raw_aggregated_temperature}
	\end{subfigure}
	\hfill
	\begin{subfigure}[b]{0.4\textwidth}
		\centering
		\includegraphics[
		width=\textwidth
		]{[Brandenburg][Manhattan][CI]Aggregated_average_temperature.png}
		\caption{}
		\label{fig:brandenburg_Manhattan_ci_aggregated_temperature}
	\end{subfigure}
	
	\caption{
		Temperatura media del aire para Brandenburgo obtenida mediante la distancia Euclidiana. Las dos primeras representaciones muestran los resultados para los datos sin agregar y sus respectivos intervalos de confianza, mientras que las dos representaciones inferiores muestran los resultados agregados para los mismos casos.
	}
	\label{fig:brandenburg_Manhattan_average_temperature}
	
	
\end{figure}




\begin{figure}[htbp]
	\centering
	
	
	% ============================================================
	% Primera fila: Raw
	% ============================================================
	\begin{subfigure}[b]{0.8\textwidth}
		\includegraphics[
		width=\textwidth
		]{[Brandenburg][Euclidean][Raw]relative_humidity.png}
		\caption{}
		\label{fig:brandenburg_euclidean_raw_yearly_relative_humidity}
	\end{subfigure}
	\vspace{2mm}
	
	% ============================================================
	% Segunda fila: CI
	% ============================================================
	\begin{subfigure}[b]{0.8\textwidth}
		\includegraphics[
		width=\textwidth
		]{[Brandenburg][Euclidean][CI]relative_humidity.png}
		\caption{}
		\label{fig:brandenburg_euclidean_ci_yearly_relative_humidity}
	\end{subfigure}
	\vspace{2mm}
	
	% ============================================================
	% Tercera fila: datos agregados
	% ============================================================
	
	\begin{subfigure}[b]{0.4\textwidth}
		\centering
		\includegraphics[
		width=\textwidth
		]{[Brandenburg][Euclidean][Raw]Aggregated_relative_humidity.png}
		\caption{}
		\label{fig:brandenburg_euclidean_raw_aggregated_relative_humidity}
	\end{subfigure}
	\hfill
	\begin{subfigure}[b]{0.4\textwidth}
		\centering
		\includegraphics[
		width=\textwidth
		]{[Brandenburg][Euclidean][CI]Aggregated_relative_humidity.png}
		\caption{}
		\label{fig:brandenburg_euclidean_ci_aggregated_relative_humidity}
	\end{subfigure}
	
	\caption{
		Temperatura media del aire para Brandenburgo obtenida mediante la distancia Euclidiana. Las dos primeras representaciones muestran los resultados para los datos sin agregar y sus respectivos intervalos de confianza, mientras que las dos representaciones inferiores muestran los resultados agregados para los mismos casos.
	}
	\label{fig:brandenburg_euclidean_relative_humidity}
	
	
\end{figure}



\begin{figure}[htbp]
	\centering
	
	
	% ============================================================
	% Primera fila: Raw
	% ============================================================
	\begin{subfigure}[b]{0.8\textwidth}
		\includegraphics[
		width=\textwidth
		]{[Brandenburg][Manhattan][Raw]relative_humidity.png}
		\caption{}
		\label{fig:brandenburg_Manhattan_raw_yearly_relative_humidity}
	\end{subfigure}
	\vspace{2mm}
	
	% ============================================================
	% Segunda fila: CI
	% ============================================================
	\begin{subfigure}[b]{0.8\textwidth}
		\includegraphics[
		width=\textwidth
		]{[Brandenburg][Manhattan][CI]relative_humidity.png}
		\caption{}
		\label{fig:brandenburg_Manhattan_ci_yearly_relative_humidity}
	\end{subfigure}
	\vspace{2mm}
	
	% ============================================================
	% Tercera fila: datos agregados
	% ============================================================
	
	\begin{subfigure}[b]{0.4\textwidth}
		\centering
		\includegraphics[
		width=\textwidth
		]{[Brandenburg][Manhattan][Raw]Aggregated_relative_humidity.png}
		\caption{}
		\label{fig:brandenburg_Manhattan_raw_aggregated_relative_humidity}
	\end{subfigure}
	\hfill
	\begin{subfigure}[b]{0.4\textwidth}
		\centering
		\includegraphics[
		width=\textwidth
		]{[Brandenburg][Manhattan][CI]Aggregated_relative_humidity.png}
		\caption{}
		\label{fig:brandenburg_Manhattan_ci_aggregated_relative_humidity}
	\end{subfigure}
	
	\caption{
		Temperatura media del aire para Brandenburgo obtenida mediante la distancia Euclidiana. Las dos primeras representaciones muestran los resultados para los datos sin agregar y sus respectivos intervalos de confianza, mientras que las dos representaciones inferiores muestran los resultados agregados para los mismos casos.
	}
	\label{fig:brandenburg_Manhattan_relative_humidity}
	
	
\end{figure}





\begin{figure}[htbp]
	\centering
	
	
	% ============================================================
	% Primera fila: Raw
	% ============================================================
	\begin{subfigure}[b]{0.8\textwidth}
		\includegraphics[
		width=\textwidth
		]{[Brandenburg][Euclidean][Raw]soil_moisture_0-10cm.png}
		\caption{}
		\label{fig:brandenburg_euclidean_raw_yearly_soil_moisture_0-10cm}
	\end{subfigure}
	\vspace{2mm}
	
	% ============================================================
	% Segunda fila: CI
	% ============================================================
	\begin{subfigure}[b]{0.8\textwidth}
		\includegraphics[
		width=\textwidth
		]{[Brandenburg][Euclidean][CI]soil_moisture_0-10cm.png}
		\caption{}
		\label{fig:brandenburg_euclidean_ci_yearly_soil_moisture_0-10cm}
	\end{subfigure}
	\vspace{2mm}
	
	% ============================================================
	% Tercera fila: datos agregados
	% ============================================================
	
	\begin{subfigure}[b]{0.4\textwidth}
		\centering
		\includegraphics[
		width=\textwidth
		]{[Brandenburg][Euclidean][Raw]Aggregated_soil_moisture_0-10cm.png}
		\caption{}
		\label{fig:brandenburg_euclidean_raw_aggregated_soil_moisture_0-10cm}
	\end{subfigure}
	\hfill
	\begin{subfigure}[b]{0.4\textwidth}
		\centering
		\includegraphics[
		width=\textwidth
		]{[Brandenburg][Euclidean][CI]Aggregated_soil_moisture_0-10cm.png}
		\caption{}
		\label{fig:brandenburg_euclidean_ci_aggregated_soil_moisture_0-10cm}
	\end{subfigure}
	
	\caption{
		Temperatura media del aire para Brandenburgo obtenida mediante la distancia Euclidiana. Las dos primeras representaciones muestran los resultados para los datos sin agregar y sus respectivos intervalos de confianza, mientras que las dos representaciones inferiores muestran los resultados agregados para los mismos casos.
	}
	\label{fig:brandenburg_euclidean_soil_moisture_0-10cm}
	
	
\end{figure}



\begin{figure}[htbp]
	\centering
	
	
	% ============================================================
	% Primera fila: Raw
	% ============================================================
	\begin{subfigure}[b]{0.8\textwidth}
		\includegraphics[
		width=\textwidth
		]{[Brandenburg][Manhattan][Raw]soil_moisture_0-10cm.png}
		\caption{}
		\label{fig:brandenburg_Manhattan_raw_yearly_soil_moisture_0-10cm}
	\end{subfigure}
	\vspace{2mm}
	
	% ============================================================
	% Segunda fila: CI
	% ============================================================
	\begin{subfigure}[b]{0.8\textwidth}
		\includegraphics[
		width=\textwidth
		]{[Brandenburg][Manhattan][CI]soil_moisture_0-10cm.png}
		\caption{}
		\label{fig:brandenburg_Manhattan_ci_yearly_soil_moisture_0-10cm}
	\end{subfigure}
	\vspace{2mm}
	
	% ============================================================
	% Tercera fila: datos agregados
	% ============================================================
	
	\begin{subfigure}[b]{0.4\textwidth}
		\centering
		\includegraphics[
		width=\textwidth
		]{[Brandenburg][Manhattan][Raw]Aggregated_soil_moisture_0-10cm.png}
		\caption{}
		\label{fig:brandenburg_Manhattan_raw_aggregated_soil_moisture_0-10cm}
	\end{subfigure}
	\hfill
	\begin{subfigure}[b]{0.4\textwidth}
		\centering
		\includegraphics[
		width=\textwidth
		]{[Brandenburg][Manhattan][CI]Aggregated_soil_moisture_0-10cm.png}
		\caption{}
		\label{fig:brandenburg_Manhattan_ci_aggregated_soil_moisture_0-10cm}
	\end{subfigure}
	
	\caption{
		Temperatura media del aire para Brandenburgo obtenida mediante la distancia Euclidiana. Las dos primeras representaciones muestran los resultados para los datos sin agregar y sus respectivos intervalos de confianza, mientras que las dos representaciones inferiores muestran los resultados agregados para los mismos casos.
	}
	\label{fig:brandenburg_Manhattan_soil_moisture_0-10cm}
	
	
\end{figure}


\begin{figure}[htbp]
	\centering
	
	
	% ============================================================
	% Primera fila: Raw
	% ============================================================
	\begin{subfigure}[b]{0.8\textwidth}
		\includegraphics[
		width=\textwidth
		]{[Brandenburg][Euclidean][Raw]soil_temperature_5cm.png}
		\caption{}
		\label{fig:brandenburg_euclidean_raw_yearly_soil_temperature_5cm}
	\end{subfigure}
	\vspace{2mm}
	
	% ============================================================
	% Segunda fila: CI
	% ============================================================
	\begin{subfigure}[b]{0.8\textwidth}
		\includegraphics[
		width=\textwidth
		]{[Brandenburg][Euclidean][CI]soil_temperature_5cm.png}
		\caption{}
		\label{fig:brandenburg_euclidean_ci_yearly_soil_temperature_5cm}
	\end{subfigure}
	\vspace{2mm}
	
	% ============================================================
	% Tercera fila: datos agregados
	% ============================================================
	
	\begin{subfigure}[b]{0.4\textwidth}
		\centering
		\includegraphics[
		width=\textwidth
		]{[Brandenburg][Euclidean][Raw]Aggregated_soil_temperature_5cm.png}
		\caption{}
		\label{fig:brandenburg_euclidean_raw_aggregated_soil_temperature_5cm}
	\end{subfigure}
	\hfill
	\begin{subfigure}[b]{0.4\textwidth}
		\centering
		\includegraphics[
		width=\textwidth
		]{[Brandenburg][Euclidean][CI]Aggregated_soil_temperature_5cm.png}
		\caption{}
		\label{fig:brandenburg_euclidean_ci_aggregated_soil_temperature_5cm}
	\end{subfigure}
	
	\caption{
		Temperatura media del aire para Brandenburgo obtenida mediante la distancia Euclidiana. Las dos primeras representaciones muestran los resultados para los datos sin agregar y sus respectivos intervalos de confianza, mientras que las dos representaciones inferiores muestran los resultados agregados para los mismos casos.
	}
	\label{fig:brandenburg_euclidean_soil_temperature_5cm}
	
	
\end{figure}



\begin{figure}[htbp]
	\centering
	
	
	% ============================================================
	% Primera fila: Raw
	% ============================================================
	\begin{subfigure}[b]{0.8\textwidth}
		\includegraphics[
		width=\textwidth
		]{[Brandenburg][Manhattan][Raw]soil_temperature_5cm.png}
		\caption{}
		\label{fig:brandenburg_Manhattan_raw_yearly_soil_temperature_5cm}
	\end{subfigure}
	\vspace{2mm}
	
	% ============================================================
	% Segunda fila: CI
	% ============================================================
	\begin{subfigure}[b]{0.8\textwidth}
		\includegraphics[
		width=\textwidth
		]{[Brandenburg][Manhattan][CI]soil_temperature_5cm.png}
		\caption{}
		\label{fig:brandenburg_Manhattan_ci_yearly_soil_temperature_5cm}
	\end{subfigure}
	\vspace{2mm}
	
	% ============================================================
	% Tercera fila: datos agregados
	% ============================================================
	
	\begin{subfigure}[b]{0.4\textwidth}
		\centering
		\includegraphics[
		width=\textwidth
		]{[Brandenburg][Manhattan][Raw]Aggregated_soil_temperature_5cm.png}
		\caption{}
		\label{fig:brandenburg_Manhattan_raw_aggregated_soil_temperature_5cm}
	\end{subfigure}
	\hfill
	\begin{subfigure}[b]{0.4\textwidth}
		\centering
		\includegraphics[
		width=\textwidth
		]{[Brandenburg][Manhattan][CI]Aggregated_soil_temperature_5cm.png}
		\caption{}
		\label{fig:brandenburg_Manhattan_ci_aggregated_soil_temperature_5cm}
	\end{subfigure}
	
	\caption{
		Temperatura media del aire para Brandenburgo obtenida mediante la distancia Euclidiana. Las dos primeras representaciones muestran los resultados para los datos sin agregar y sus respectivos intervalos de confianza, mientras que las dos representaciones inferiores muestran los resultados agregados para los mismos casos.
	}
	\label{fig:brandenburg_Manhattan_soil_temperature_5cm}
	
	
\end{figure}





%################################
%								#
%	        SAXONY  			#
%								#
%################################



\begin{figure}[htbp]
	\centering
	
	
	% ============================================================
	% Primera fila: Raw
	% ============================================================
	\begin{subfigure}[b]{0.8\textwidth}
		\includegraphics[
		width=\textwidth
		]{[Saxony][Euclidean][Raw]average_temperature.png}
		\caption{}
		\label{fig:saxony_euclidean_raw_yearly_temperature}
	\end{subfigure}
	\vspace{2mm}
	
	% ============================================================
	% Segunda fila: CI
	% ============================================================
	\begin{subfigure}[b]{0.8\textwidth}
		\includegraphics[
		width=\textwidth
		]{[Saxony][Euclidean][CI]average_temperature.png}
		\caption{}
		\label{fig:saxony_euclidean_ci_yearly_temperature}
	\end{subfigure}
	\vspace{2mm}
	
	% ============================================================
	% Tercera fila: datos agregados
	% ============================================================
	
	\begin{subfigure}[b]{0.4\textwidth}
		\centering
		\includegraphics[
		width=\textwidth
		]{[Saxony][Euclidean][Raw]Aggregated_average_temperature.png}
		\caption{}
		\label{fig:saxony_euclidean_raw_aggregated_temperature}
	\end{subfigure}
	\hfill
	\begin{subfigure}[b]{0.4\textwidth}
		\centering
		\includegraphics[
		width=\textwidth
		]{[Saxony][Euclidean][CI]Aggregated_average_temperature.png}
		\caption{}
		\label{fig:saxony_euclidean_ci_aggregated_temperature}
	\end{subfigure}
	
	\caption{
		Temperatura media del aire para Saxonyo obtenida mediante la distancia Euclidiana. Las dos primeras representaciones muestran los resultados para los datos sin agregar y sus respectivos intervalos de confianza, mientras que las dos representaciones inferiores muestran los resultados agregados para los mismos casos.
	}
	\label{fig:saxony_euclidean_average_temperature}
	
	
\end{figure}



\begin{figure}[htbp]
	\centering
	
	
	% ============================================================
	% Primera fila: Raw
	% ============================================================
	\begin{subfigure}[b]{0.8\textwidth}
		\includegraphics[
		width=\textwidth
		]{[Saxony][Manhattan][Raw]average_temperature.png}
		\caption{}
		\label{fig:saxony_Manhattan_raw_yearly_temperature}
	\end{subfigure}
	\vspace{2mm}
	
	% ============================================================
	% Segunda fila: CI
	% ============================================================
	\begin{subfigure}[b]{0.8\textwidth}
		\includegraphics[
		width=\textwidth
		]{[Saxony][Manhattan][CI]average_temperature.png}
		\caption{}
		\label{fig:saxony_Manhattan_ci_yearly_temperature}
	\end{subfigure}
	\vspace{2mm}
	
	% ============================================================
	% Tercera fila: datos agregados
	% ============================================================
	
	\begin{subfigure}[b]{0.4\textwidth}
		\centering
		\includegraphics[
		width=\textwidth
		]{[Saxony][Manhattan][Raw]Aggregated_average_temperature.png}
		\caption{}
		\label{fig:saxony_Manhattan_raw_aggregated_temperature}
	\end{subfigure}
	\hfill
	\begin{subfigure}[b]{0.4\textwidth}
		\centering
		\includegraphics[
		width=\textwidth
		]{[Saxony][Manhattan][CI]Aggregated_average_temperature.png}
		\caption{}
		\label{fig:saxony_Manhattan_ci_aggregated_temperature}
	\end{subfigure}
	
	\caption{
		Temperatura media del aire para Saxonyo obtenida mediante la distancia Euclidiana. Las dos primeras representaciones muestran los resultados para los datos sin agregar y sus respectivos intervalos de confianza, mientras que las dos representaciones inferiores muestran los resultados agregados para los mismos casos.
	}
	\label{fig:saxony_Manhattan_average_temperature}
	
	
\end{figure}




\begin{figure}[htbp]
	\centering
	
	
	% ============================================================
	% Primera fila: Raw
	% ============================================================
	\begin{subfigure}[b]{0.8\textwidth}
		\includegraphics[
		width=\textwidth
		]{[Saxony][Euclidean][Raw]relative_humidity.png}
		\caption{}
		\label{fig:saxony_euclidean_raw_yearly_relative_humidity}
	\end{subfigure}
	\vspace{2mm}
	
	% ============================================================
	% Segunda fila: CI
	% ============================================================
	\begin{subfigure}[b]{0.8\textwidth}
		\includegraphics[
		width=\textwidth
		]{[Saxony][Euclidean][CI]relative_humidity.png}
		\caption{}
		\label{fig:saxony_euclidean_ci_yearly_relative_humidity}
	\end{subfigure}
	\vspace{2mm}
	
	% ============================================================
	% Tercera fila: datos agregados
	% ============================================================
	
	\begin{subfigure}[b]{0.4\textwidth}
		\centering
		\includegraphics[
		width=\textwidth
		]{[Saxony][Euclidean][Raw]Aggregated_relative_humidity.png}
		\caption{}
		\label{fig:saxony_euclidean_raw_aggregated_relative_humidity}
	\end{subfigure}
	\hfill
	\begin{subfigure}[b]{0.4\textwidth}
		\centering
		\includegraphics[
		width=\textwidth
		]{[Saxony][Euclidean][CI]Aggregated_relative_humidity.png}
		\caption{}
		\label{fig:saxony_euclidean_ci_aggregated_relative_humidity}
	\end{subfigure}
	
	\caption{
		Temperatura media del aire para Saxonyo obtenida mediante la distancia Euclidiana. Las dos primeras representaciones muestran los resultados para los datos sin agregar y sus respectivos intervalos de confianza, mientras que las dos representaciones inferiores muestran los resultados agregados para los mismos casos.
	}
	\label{fig:saxony_euclidean_relative_humidity}
	
	
\end{figure}



\begin{figure}[htbp]
	\centering
	
	
	% ============================================================
	% Primera fila: Raw
	% ============================================================
	\begin{subfigure}[b]{0.8\textwidth}
		\includegraphics[
		width=\textwidth
		]{[Saxony][Manhattan][Raw]relative_humidity.png}
		\caption{}
		\label{fig:saxony_Manhattan_raw_yearly_relative_humidity}
	\end{subfigure}
	\vspace{2mm}
	
	% ============================================================
	% Segunda fila: CI
	% ============================================================
	\begin{subfigure}[b]{0.8\textwidth}
		\includegraphics[
		width=\textwidth
		]{[Saxony][Manhattan][CI]relative_humidity.png}
		\caption{}
		\label{fig:saxony_Manhattan_ci_yearly_relative_humidity}
	\end{subfigure}
	\vspace{2mm}
	
	% ============================================================
	% Tercera fila: datos agregados
	% ============================================================
	
	\begin{subfigure}[b]{0.4\textwidth}
		\centering
		\includegraphics[
		width=\textwidth
		]{[Saxony][Manhattan][Raw]Aggregated_relative_humidity.png}
		\caption{}
		\label{fig:saxony_Manhattan_raw_aggregated_relative_humidity}
	\end{subfigure}
	\hfill
	\begin{subfigure}[b]{0.4\textwidth}
		\centering
		\includegraphics[
		width=\textwidth
		]{[Saxony][Manhattan][CI]Aggregated_relative_humidity.png}
		\caption{}
		\label{fig:saxony_Manhattan_ci_aggregated_relative_humidity}
	\end{subfigure}
	
	\caption{
		Temperatura media del aire para Saxonyo obtenida mediante la distancia Euclidiana. Las dos primeras representaciones muestran los resultados para los datos sin agregar y sus respectivos intervalos de confianza, mientras que las dos representaciones inferiores muestran los resultados agregados para los mismos casos.
	}
	\label{fig:saxony_Manhattan_relative_humidity}
	
	
\end{figure}





\begin{figure}[htbp]
	\centering
	
	
	% ============================================================
	% Primera fila: Raw
	% ============================================================
	\begin{subfigure}[b]{0.8\textwidth}
		\includegraphics[
		width=\textwidth
		]{[Saxony][Euclidean][Raw]soil_moisture_0-10cm.png}
		\caption{}
		\label{fig:saxony_euclidean_raw_yearly_soil_moisture_0-10cm}
	\end{subfigure}
	\vspace{2mm}
	
	% ============================================================
	% Segunda fila: CI
	% ============================================================
	\begin{subfigure}[b]{0.8\textwidth}
		\includegraphics[
		width=\textwidth
		]{[Saxony][Euclidean][CI]soil_moisture_0-10cm.png}
		\caption{}
		\label{fig:saxony_euclidean_ci_yearly_soil_moisture_0-10cm}
	\end{subfigure}
	\vspace{2mm}
	
	% ============================================================
	% Tercera fila: datos agregados
	% ============================================================
	
	\begin{subfigure}[b]{0.4\textwidth}
		\centering
		\includegraphics[
		width=\textwidth
		]{[Saxony][Euclidean][Raw]Aggregated_soil_moisture_0-10cm.png}
		\caption{}
		\label{fig:saxony_euclidean_raw_aggregated_soil_moisture_0-10cm}
	\end{subfigure}
	\hfill
	\begin{subfigure}[b]{0.4\textwidth}
		\centering
		\includegraphics[
		width=\textwidth
		]{[Saxony][Euclidean][CI]Aggregated_soil_moisture_0-10cm.png}
		\caption{}
		\label{fig:saxony_euclidean_ci_aggregated_soil_moisture_0-10cm}
	\end{subfigure}
	
	\caption{
		Temperatura media del aire para Saxonyo obtenida mediante la distancia Euclidiana. Las dos primeras representaciones muestran los resultados para los datos sin agregar y sus respectivos intervalos de confianza, mientras que las dos representaciones inferiores muestran los resultados agregados para los mismos casos.
	}
	\label{fig:saxony_euclidean_soil_moisture_0-10cm}
	
	
\end{figure}



\begin{figure}[htbp]
	\centering
	
	
	% ============================================================
	% Primera fila: Raw
	% ============================================================
	\begin{subfigure}[b]{0.8\textwidth}
		\includegraphics[
		width=\textwidth
		]{[Saxony][Manhattan][Raw]soil_moisture_0-10cm.png}
		\caption{}
		\label{fig:saxony_Manhattan_raw_yearly_soil_moisture_0-10cm}
	\end{subfigure}
	\vspace{2mm}
	
	% ============================================================
	% Segunda fila: CI
	% ============================================================
	\begin{subfigure}[b]{0.8\textwidth}
		\includegraphics[
		width=\textwidth
		]{[Saxony][Manhattan][CI]soil_moisture_0-10cm.png}
		\caption{}
		\label{fig:saxony_Manhattan_ci_yearly_soil_moisture_0-10cm}
	\end{subfigure}
	\vspace{2mm}
	
	% ============================================================
	% Tercera fila: datos agregados
	% ============================================================
	
	\begin{subfigure}[b]{0.4\textwidth}
		\centering
		\includegraphics[
		width=\textwidth
		]{[Saxony][Manhattan][Raw]Aggregated_soil_moisture_0-10cm.png}
		\caption{}
		\label{fig:saxony_Manhattan_raw_aggregated_soil_moisture_0-10cm}
	\end{subfigure}
	\hfill
	\begin{subfigure}[b]{0.4\textwidth}
		\centering
		\includegraphics[
		width=\textwidth
		]{[Saxony][Manhattan][CI]Aggregated_soil_moisture_0-10cm.png}
		\caption{}
		\label{fig:saxony_Manhattan_ci_aggregated_soil_moisture_0-10cm}
	\end{subfigure}
	
	\caption{
		Temperatura media del aire para Saxonyo obtenida mediante la distancia Euclidiana. Las dos primeras representaciones muestran los resultados para los datos sin agregar y sus respectivos intervalos de confianza, mientras que las dos representaciones inferiores muestran los resultados agregados para los mismos casos.
	}
	\label{fig:saxony_Manhattan_soil_moisture_0-10cm}
	
	
\end{figure}


\begin{figure}[htbp]
	\centering
	
	
	% ============================================================
	% Primera fila: Raw
	% ============================================================
	\begin{subfigure}[b]{0.8\textwidth}
		\includegraphics[
		width=\textwidth
		]{[Saxony][Euclidean][Raw]soil_temperature_5cm.png}
		\caption{}
		\label{fig:saxony_euclidean_raw_yearly_soil_temperature_5cm}
	\end{subfigure}
	\vspace{2mm}
	
	% ============================================================
	% Segunda fila: CI
	% ============================================================
	\begin{subfigure}[b]{0.8\textwidth}
		\includegraphics[
		width=\textwidth
		]{[Saxony][Euclidean][CI]soil_temperature_5cm.png}
		\caption{}
		\label{fig:saxony_euclidean_ci_yearly_soil_temperature_5cm}
	\end{subfigure}
	\vspace{2mm}
	
	% ============================================================
	% Tercera fila: datos agregados
	% ============================================================
	
	\begin{subfigure}[b]{0.4\textwidth}
		\centering
		\includegraphics[
		width=\textwidth
		]{[Saxony][Euclidean][Raw]Aggregated_soil_temperature_5cm.png}
		\caption{}
		\label{fig:saxony_euclidean_raw_aggregated_soil_temperature_5cm}
	\end{subfigure}
	\hfill
	\begin{subfigure}[b]{0.4\textwidth}
		\centering
		\includegraphics[
		width=\textwidth
		]{[Saxony][Euclidean][CI]Aggregated_soil_temperature_5cm.png}
		\caption{}
		\label{fig:saxony_euclidean_ci_aggregated_soil_temperature_5cm}
	\end{subfigure}
	
	\caption{
		Temperatura media del aire para Saxonyo obtenida mediante la distancia Euclidiana. Las dos primeras representaciones muestran los resultados para los datos sin agregar y sus respectivos intervalos de confianza, mientras que las dos representaciones inferiores muestran los resultados agregados para los mismos casos.
	}
	\label{fig:saxony_euclidean_soil_temperature_5cm}
	
	
\end{figure}



\begin{figure}[htbp]
	\centering
	
	
	% ============================================================
	% Primera fila: Raw
	% ============================================================
	\begin{subfigure}[b]{0.8\textwidth}
		\includegraphics[
		width=\textwidth
		]{[Saxony][Manhattan][Raw]soil_temperature_5cm.png}
		\caption{}
		\label{fig:saxony_Manhattan_raw_yearly_soil_temperature_5cm}
	\end{subfigure}
	\vspace{2mm}
	
	% ============================================================
	% Segunda fila: CI
	% ============================================================
	\begin{subfigure}[b]{0.8\textwidth}
		\includegraphics[
		width=\textwidth
		]{[Saxony][Manhattan][CI]soil_temperature_5cm.png}
		\caption{}
		\label{fig:saxony_Manhattan_ci_yearly_soil_temperature_5cm}
	\end{subfigure}
	\vspace{2mm}
	
	% ============================================================
	% Tercera fila: datos agregados
	% ============================================================
	
	\begin{subfigure}[b]{0.4\textwidth}
		\centering
		\includegraphics[
		width=\textwidth
		]{[Saxony][Manhattan][Raw]Aggregated_soil_temperature_5cm.png}
		\caption{}
		\label{fig:saxony_Manhattan_raw_aggregated_soil_temperature_5cm}
	\end{subfigure}
	\hfill
	\begin{subfigure}[b]{0.4\textwidth}
		\centering
		\includegraphics[
		width=\textwidth
		]{[Saxony][Manhattan][CI]Aggregated_soil_temperature_5cm.png}
		\caption{}
		\label{fig:saxony_Manhattan_ci_aggregated_soil_temperature_5cm}
	\end{subfigure}
	
	\caption{
		Temperatura media del aire para Saxonyo obtenida mediante la distancia Euclidiana. Las dos primeras representaciones muestran los resultados para los datos sin agregar y sus respectivos intervalos de confianza, mientras que las dos representaciones inferiores muestran los resultados agregados para los mismos casos.
	}
	\label{fig:saxony_Manhattan_soil_temperature_5cm}
	
	
\end{figure}





%################################
%								#
%	        BAVARIA  			#
%								#
%################################



\begin{figure}[htbp]
	\centering
	
	
	% ============================================================
	% Primera fila: Raw
	% ============================================================
	\begin{subfigure}[b]{0.8\textwidth}
		\includegraphics[
		width=\textwidth
		]{[Bavaria][Euclidean][Raw]average_temperature.png}
		\caption{}
		\label{fig:bavaria_euclidean_raw_yearly_temperature}
	\end{subfigure}
	\vspace{2mm}
	
	% ============================================================
	% Segunda fila: CI
	% ============================================================
	\begin{subfigure}[b]{0.8\textwidth}
		\includegraphics[
		width=\textwidth
		]{[Bavaria][Euclidean][CI]average_temperature.png}
		\caption{}
		\label{fig:bavaria_euclidean_ci_yearly_temperature}
	\end{subfigure}
	\vspace{2mm}
	
	% ============================================================
	% Tercera fila: datos agregados
	% ============================================================
	
	\begin{subfigure}[b]{0.4\textwidth}
		\centering
		\includegraphics[
		width=\textwidth
		]{[Bavaria][Euclidean][Raw]Aggregated_average_temperature.png}
		\caption{}
		\label{fig:bavaria_euclidean_raw_aggregated_temperature}
	\end{subfigure}
	\hfill
	\begin{subfigure}[b]{0.4\textwidth}
		\centering
		\includegraphics[
		width=\textwidth
		]{[Bavaria][Euclidean][CI]Aggregated_average_temperature.png}
		\caption{}
		\label{fig:bavaria_euclidean_ci_aggregated_temperature}
	\end{subfigure}
	
	\caption{
		Temperatura media del aire para Bavariao obtenida mediante la distancia Euclidiana. Las dos primeras representaciones muestran los resultados para los datos sin agregar y sus respectivos intervalos de confianza, mientras que las dos representaciones inferiores muestran los resultados agregados para los mismos casos.
	}
	\label{fig:bavaria_euclidean_average_temperature}
	
	
\end{figure}



\begin{figure}[htbp]
	\centering
	
	
	% ============================================================
	% Primera fila: Raw
	% ============================================================
	\begin{subfigure}[b]{0.8\textwidth}
		\includegraphics[
		width=\textwidth
		]{[Bavaria][Manhattan][Raw]average_temperature.png}
		\caption{}
		\label{fig:bavaria_Manhattan_raw_yearly_temperature}
	\end{subfigure}
	\vspace{2mm}
	
	% ============================================================
	% Segunda fila: CI
	% ============================================================
	\begin{subfigure}[b]{0.8\textwidth}
		\includegraphics[
		width=\textwidth
		]{[Bavaria][Manhattan][CI]average_temperature.png}
		\caption{}
		\label{fig:bavaria_Manhattan_ci_yearly_temperature}
	\end{subfigure}
	\vspace{2mm}
	
	% ============================================================
	% Tercera fila: datos agregados
	% ============================================================
	
	\begin{subfigure}[b]{0.4\textwidth}
		\centering
		\includegraphics[
		width=\textwidth
		]{[Bavaria][Manhattan][Raw]Aggregated_average_temperature.png}
		\caption{}
		\label{fig:bavaria_Manhattan_raw_aggregated_temperature}
	\end{subfigure}
	\hfill
	\begin{subfigure}[b]{0.4\textwidth}
		\centering
		\includegraphics[
		width=\textwidth
		]{[Bavaria][Manhattan][CI]Aggregated_average_temperature.png}
		\caption{}
		\label{fig:bavaria_Manhattan_ci_aggregated_temperature}
	\end{subfigure}
	
	\caption{
		Temperatura media del aire para Bavariao obtenida mediante la distancia Euclidiana. Las dos primeras representaciones muestran los resultados para los datos sin agregar y sus respectivos intervalos de confianza, mientras que las dos representaciones inferiores muestran los resultados agregados para los mismos casos.
	}
	\label{fig:bavaria_Manhattan_average_temperature}
	
	
\end{figure}




\begin{figure}[htbp]
	\centering
	
	
	% ============================================================
	% Primera fila: Raw
	% ============================================================
	\begin{subfigure}[b]{0.8\textwidth}
		\includegraphics[
		width=\textwidth
		]{[Bavaria][Euclidean][Raw]relative_humidity.png}
		\caption{}
		\label{fig:bavaria_euclidean_raw_yearly_relative_humidity}
	\end{subfigure}
	\vspace{2mm}
	
	% ============================================================
	% Segunda fila: CI
	% ============================================================
	\begin{subfigure}[b]{0.8\textwidth}
		\includegraphics[
		width=\textwidth
		]{[Bavaria][Euclidean][CI]relative_humidity.png}
		\caption{}
		\label{fig:bavaria_euclidean_ci_yearly_relative_humidity}
	\end{subfigure}
	\vspace{2mm}
	
	% ============================================================
	% Tercera fila: datos agregados
	% ============================================================
	
	\begin{subfigure}[b]{0.4\textwidth}
		\centering
		\includegraphics[
		width=\textwidth
		]{[Bavaria][Euclidean][Raw]Aggregated_relative_humidity.png}
		\caption{}
		\label{fig:bavaria_euclidean_raw_aggregated_relative_humidity}
	\end{subfigure}
	\hfill
	\begin{subfigure}[b]{0.4\textwidth}
		\centering
		\includegraphics[
		width=\textwidth
		]{[Bavaria][Euclidean][CI]Aggregated_relative_humidity.png}
		\caption{}
		\label{fig:bavaria_euclidean_ci_aggregated_relative_humidity}
	\end{subfigure}
	
	\caption{
		Temperatura media del aire para Bavariao obtenida mediante la distancia Euclidiana. Las dos primeras representaciones muestran los resultados para los datos sin agregar y sus respectivos intervalos de confianza, mientras que las dos representaciones inferiores muestran los resultados agregados para los mismos casos.
	}
	\label{fig:bavaria_euclidean_relative_humidity}
	
	
\end{figure}



\begin{figure}[htbp]
	\centering
	
	
	% ============================================================
	% Primera fila: Raw
	% ============================================================
	\begin{subfigure}[b]{0.8\textwidth}
		\includegraphics[
		width=\textwidth
		]{[Bavaria][Manhattan][Raw]relative_humidity.png}
		\caption{}
		\label{fig:bavaria_Manhattan_raw_yearly_relative_humidity}
	\end{subfigure}
	\vspace{2mm}
	
	% ============================================================
	% Segunda fila: CI
	% ============================================================
	\begin{subfigure}[b]{0.8\textwidth}
		\includegraphics[
		width=\textwidth
		]{[Bavaria][Manhattan][CI]relative_humidity.png}
		\caption{}
		\label{fig:bavaria_Manhattan_ci_yearly_relative_humidity}
	\end{subfigure}
	\vspace{2mm}
	
	% ============================================================
	% Tercera fila: datos agregados
	% ============================================================
	
	\begin{subfigure}[b]{0.4\textwidth}
		\centering
		\includegraphics[
		width=\textwidth
		]{[Bavaria][Manhattan][Raw]Aggregated_relative_humidity.png}
		\caption{}
		\label{fig:bavaria_Manhattan_raw_aggregated_relative_humidity}
	\end{subfigure}
	\hfill
	\begin{subfigure}[b]{0.4\textwidth}
		\centering
		\includegraphics[
		width=\textwidth
		]{[Bavaria][Manhattan][CI]Aggregated_relative_humidity.png}
		\caption{}
		\label{fig:bavaria_Manhattan_ci_aggregated_relative_humidity}
	\end{subfigure}
	
	\caption{
		Temperatura media del aire para Bavariao obtenida mediante la distancia Euclidiana. Las dos primeras representaciones muestran los resultados para los datos sin agregar y sus respectivos intervalos de confianza, mientras que las dos representaciones inferiores muestran los resultados agregados para los mismos casos.
	}
	\label{fig:bavaria_Manhattan_relative_humidity}
	
	
\end{figure}





\begin{figure}[htbp]
	\centering
	
	
	% ============================================================
	% Primera fila: Raw
	% ============================================================
	\begin{subfigure}[b]{0.8\textwidth}
		\includegraphics[
		width=\textwidth
		]{[Bavaria][Euclidean][Raw]soil_moisture_0-10cm.png}
		\caption{}
		\label{fig:bavaria_euclidean_raw_yearly_soil_moisture_0-10cm}
	\end{subfigure}
	\vspace{2mm}
	
	% ============================================================
	% Segunda fila: CI
	% ============================================================
	\begin{subfigure}[b]{0.8\textwidth}
		\includegraphics[
		width=\textwidth
		]{[Bavaria][Euclidean][CI]soil_moisture_0-10cm.png}
		\caption{}
		\label{fig:bavaria_euclidean_ci_yearly_soil_moisture_0-10cm}
	\end{subfigure}
	\vspace{2mm}
	
	% ============================================================
	% Tercera fila: datos agregados
	% ============================================================
	
	\begin{subfigure}[b]{0.4\textwidth}
		\centering
		\includegraphics[
		width=\textwidth
		]{[Bavaria][Euclidean][Raw]Aggregated_soil_moisture_0-10cm.png}
		\caption{}
		\label{fig:bavaria_euclidean_raw_aggregated_soil_moisture_0-10cm}
	\end{subfigure}
	\hfill
	\begin{subfigure}[b]{0.4\textwidth}
		\centering
		\includegraphics[
		width=\textwidth
		]{[Bavaria][Euclidean][CI]Aggregated_soil_moisture_0-10cm.png}
		\caption{}
		\label{fig:bavaria_euclidean_ci_aggregated_soil_moisture_0-10cm}
	\end{subfigure}
	
	\caption{
		Temperatura media del aire para Bavariao obtenida mediante la distancia Euclidiana. Las dos primeras representaciones muestran los resultados para los datos sin agregar y sus respectivos intervalos de confianza, mientras que las dos representaciones inferiores muestran los resultados agregados para los mismos casos.
	}
	\label{fig:bavaria_euclidean_soil_moisture_0-10cm}
	
	
\end{figure}



\begin{figure}[htbp]
	\centering
	
	
	% ============================================================
	% Primera fila: Raw
	% ============================================================
	\begin{subfigure}[b]{0.8\textwidth}
		\includegraphics[
		width=\textwidth
		]{[Bavaria][Manhattan][Raw]soil_moisture_0-10cm.png}
		\caption{}
		\label{fig:bavaria_Manhattan_raw_yearly_soil_moisture_0-10cm}
	\end{subfigure}
	\vspace{2mm}
	
	% ============================================================
	% Segunda fila: CI
	% ============================================================
	\begin{subfigure}[b]{0.8\textwidth}
		\includegraphics[
		width=\textwidth
		]{[Bavaria][Manhattan][CI]soil_moisture_0-10cm.png}
		\caption{}
		\label{fig:bavaria_Manhattan_ci_yearly_soil_moisture_0-10cm}
	\end{subfigure}
	\vspace{2mm}
	
	% ============================================================
	% Tercera fila: datos agregados
	% ============================================================
	
	\begin{subfigure}[b]{0.4\textwidth}
		\centering
		\includegraphics[
		width=\textwidth
		]{[Bavaria][Manhattan][Raw]Aggregated_soil_moisture_0-10cm.png}
		\caption{}
		\label{fig:bavaria_Manhattan_raw_aggregated_soil_moisture_0-10cm}
	\end{subfigure}
	\hfill
	\begin{subfigure}[b]{0.4\textwidth}
		\centering
		\includegraphics[
		width=\textwidth
		]{[Bavaria][Manhattan][CI]Aggregated_soil_moisture_0-10cm.png}
		\caption{}
		\label{fig:bavaria_Manhattan_ci_aggregated_soil_moisture_0-10cm}
	\end{subfigure}
	
	\caption{
		Temperatura media del aire para Bavariao obtenida mediante la distancia Euclidiana. Las dos primeras representaciones muestran los resultados para los datos sin agregar y sus respectivos intervalos de confianza, mientras que las dos representaciones inferiores muestran los resultados agregados para los mismos casos.
	}
	\label{fig:bavaria_Manhattan_soil_moisture_0-10cm}
	
	
\end{figure}


\begin{figure}[htbp]
	\centering
	
	
	% ============================================================
	% Primera fila: Raw
	% ============================================================
	\begin{subfigure}[b]{0.8\textwidth}
		\includegraphics[
		width=\textwidth
		]{[Bavaria][Euclidean][Raw]soil_temperature_5cm.png}
		\caption{}
		\label{fig:bavaria_euclidean_raw_yearly_soil_temperature_5cm}
	\end{subfigure}
	\vspace{2mm}
	
	% ============================================================
	% Segunda fila: CI
	% ============================================================
	\begin{subfigure}[b]{0.8\textwidth}
		\includegraphics[
		width=\textwidth
		]{[Bavaria][Euclidean][CI]soil_temperature_5cm.png}
		\caption{}
		\label{fig:bavaria_euclidean_ci_yearly_soil_temperature_5cm}
	\end{subfigure}
	\vspace{2mm}
	
	% ============================================================
	% Tercera fila: datos agregados
	% ============================================================
	
	\begin{subfigure}[b]{0.4\textwidth}
		\centering
		\includegraphics[
		width=\textwidth
		]{[Bavaria][Euclidean][Raw]Aggregated_soil_temperature_5cm.png}
		\caption{}
		\label{fig:bavaria_euclidean_raw_aggregated_soil_temperature_5cm}
	\end{subfigure}
	\hfill
	\begin{subfigure}[b]{0.4\textwidth}
		\centering
		\includegraphics[
		width=\textwidth
		]{[Bavaria][Euclidean][CI]Aggregated_soil_temperature_5cm.png}
		\caption{}
		\label{fig:bavaria_euclidean_ci_aggregated_soil_temperature_5cm}
	\end{subfigure}
	
	\caption{
		Temperatura media del aire para Bavariao obtenida mediante la distancia Euclidiana. Las dos primeras representaciones muestran los resultados para los datos sin agregar y sus respectivos intervalos de confianza, mientras que las dos representaciones inferiores muestran los resultados agregados para los mismos casos.
	}
	\label{fig:bavaria_euclidean_soil_temperature_5cm}
	
	
\end{figure}



\begin{figure}[htbp]
	\centering
	
	
	% ============================================================
	% Primera fila: Raw
	% ============================================================
	\begin{subfigure}[b]{0.8\textwidth}
		\includegraphics[
		width=\textwidth
		]{[Bavaria][Manhattan][Raw]soil_temperature_5cm.png}
		\caption{}
		\label{fig:bavaria_Manhattan_raw_yearly_soil_temperature_5cm}
	\end{subfigure}
	\vspace{2mm}
	
	% ============================================================
	% Segunda fila: CI
	% ============================================================
	\begin{subfigure}[b]{0.8\textwidth}
		\includegraphics[
		width=\textwidth
		]{[Bavaria][Manhattan][CI]soil_temperature_5cm.png}
		\caption{}
		\label{fig:bavaria_Manhattan_ci_yearly_soil_temperature_5cm}
	\end{subfigure}
	\vspace{2mm}
	
	% ============================================================
	% Tercera fila: datos agregados
	% ============================================================
	
	\begin{subfigure}[b]{0.4\textwidth}
		\centering
		\includegraphics[
		width=\textwidth
		]{[Bavaria][Manhattan][Raw]Aggregated_soil_temperature_5cm.png}
		\caption{}
		\label{fig:bavaria_Manhattan_raw_aggregated_soil_temperature_5cm}
	\end{subfigure}
	\hfill
	\begin{subfigure}[b]{0.4\textwidth}
		\centering
		\includegraphics[
		width=\textwidth
		]{[Bavaria][Manhattan][CI]Aggregated_soil_temperature_5cm.png}
		\caption{}
		\label{fig:bavaria_Manhattan_ci_aggregated_soil_temperature_5cm}
	\end{subfigure}
	
	\caption{
		Temperatura media del aire para Bavariao obtenida mediante la distancia Euclidiana. Las dos primeras representaciones muestran los resultados para los datos sin agregar y sus respectivos intervalos de confianza, mientras que las dos representaciones inferiores muestran los resultados agregados para los mismos casos.
	}
	\label{fig:bavaria_Manhattan_soil_temperature_5cm}
	
	
\end{figure}




\fi




De acuerdo con \cite{Flores_2025}, la región agrícola empleada se caracteriza por una elevada heterogeneidad topográfica, climática y edáfica. En particular, Baviera, Brandemburgo y Sajonia presentan variaciones en humedad, altitud, temperatura y características del suelo, factores que influyen directamente en las prácticas agrícolas y los calendarios de siembra. Esta diversidad de condiciones ambientales permitió someter la metodología propuesta a distintos escenarios espaciales y verifica su capacidad para identificar variaciones en la aptitud de las fechas de siembra y, con ello, evaluar parcialmente su comportamiento bajo condiciones ambientales heterogéneas.\\


% Please add the following required packages to your document preamble:
% \usepackage{multirow}
\begin{table}[]
	\centering
	\begin{tabular}{cclll}
		\hline
		\multicolumn{5}{l}{\textbf{Brandenburg}}                                                                                                                                                                                                                                                                                                                                 \\ \hline
		\multicolumn{1}{c|}{\multirow{3}{*}{\textbf{Semester}}} & \multicolumn{4}{c}{\textbf{Sowing windows}}                                                                                                                                                                                                                                                                    \\ \cline{2-5} 
		\multicolumn{1}{c|}{}                                   & \multicolumn{2}{c|}{\textbf{Euclidean}}                                                                                                                                                                            & \multicolumn{2}{c}{\textbf{Manhattan}}                                                    \\ \cline{2-5} 
		\multicolumn{1}{c|}{}                                   & \multicolumn{1}{c|}{\textbf{Raw prediction}}                                                    & \multicolumn{1}{c|}{\textbf{CI prediction}}                                                                      & \multicolumn{1}{c|}{\textbf{Raw prediction}} & \multicolumn{1}{c}{\textbf{CI prediction}} \\ \hline
		\multicolumn{1}{c|}{\textbf{1}}                         & \multicolumn{1}{l|}{\begin{tabular}[c]{@{}l@{}}55--108, 123--125,\\ 137, 139--142\end{tabular}} & \multicolumn{1}{l|}{\begin{tabular}[c]{@{}l@{}}55--108, 115, 117,\\ 121--127, 130--131,\\ 135--147\end{tabular}} & \multicolumn{1}{l|}{69--106}                 & 64--65, 67--107                            \\ \hline
		\multicolumn{1}{c|}{\textbf{2}}                         & \multicolumn{1}{l|}{271--310}                                                                   & \multicolumn{1}{l|}{271--310, 312, 314}                                                                          & \multicolumn{1}{l|}{273--304}                & 271--305                                   \\ \hline
		\multicolumn{1}{l}{}                                    & \multicolumn{1}{l}{}                                                                            &                                                                                                                  &                                              &                                            \\ \hline
		\multicolumn{5}{l}{\textbf{Saxony}}                                                                                                                                                                                                                                                                                                                                      \\ \hline
		\multicolumn{1}{c|}{\multirow{3}{*}{\textbf{Semester}}} & \multicolumn{4}{c}{\textbf{Sowing windows}}                                                                                                                                                                                                                                                                    \\ \cline{2-5} 
		\multicolumn{1}{c|}{}                                   & \multicolumn{2}{c|}{\textbf{Euclidean}}                                                                                                                                                                            & \multicolumn{2}{c}{\textbf{Manhattan}}                                                    \\ \cline{2-5} 
		\multicolumn{1}{c|}{}                                   & \multicolumn{1}{c|}{\textbf{Raw prediction}}                                                    & \multicolumn{1}{c|}{\textbf{CI prediction}}                                                                      & \multicolumn{1}{c|}{\textbf{Raw prediction}} & \multicolumn{1}{c}{\textbf{CI prediction}} \\ \hline
		\multicolumn{1}{c|}{\textbf{1}}                         & \multicolumn{1}{l|}{\begin{tabular}[c]{@{}l@{}}93, 96, 99,\\ 102--143, 145\end{tabular}}        & \multicolumn{1}{l|}{90, 92--145, 147}                                                                            & \multicolumn{1}{l|}{90--137, 139}            & 90--140, 142                               \\ \hline
		\multicolumn{1}{c|}{\textbf{2}}                         & \multicolumn{1}{l|}{262--274, 276}                                                              & \multicolumn{1}{l|}{260, 262--276}                                                                               & \multicolumn{1}{l|}{262--276}                & 261--276                                   \\ \hline
		\multicolumn{1}{l}{}                                    & \multicolumn{1}{l}{}                                                                            &                                                                                                                  &                                              &                                            \\ \hline
		\multicolumn{5}{l}{\textbf{Bavaria}}                                                                                                                                                                                                                                                                                                                                     \\ \hline
		\multicolumn{1}{c|}{\multirow{3}{*}{\textbf{Semester}}} & \multicolumn{4}{c}{\textbf{Sowing windows}}                                                                                                                                                                                                                                                                    \\ \cline{2-5} 
		\multicolumn{1}{c|}{}                                   & \multicolumn{2}{c|}{\textbf{Euclidean}}                                                                                                                                                                            & \multicolumn{2}{c}{\textbf{Manhattan}}                                                    \\ \cline{2-5} 
		\multicolumn{1}{c|}{}                                   & \multicolumn{1}{c|}{\textbf{Raw prediction}}                                                    & \multicolumn{1}{c|}{\textbf{CI prediction}}                                                                      & \multicolumn{1}{c|}{\textbf{Raw prediction}} & \multicolumn{1}{c}{\textbf{CI prediction}} \\ \hline
		\multicolumn{1}{c|}{\textbf{1}}                         & \multicolumn{1}{l|}{99--119}                                                                    & \multicolumn{1}{l|}{96--119}                                                                                     & \multicolumn{1}{l|}{95, 97--119}             & 91, 93--122                                \\ \hline
		\multicolumn{1}{c|}{\textbf{2}}                         & \multicolumn{1}{l|}{246--264}                                                                   & \multicolumn{1}{l|}{245--265}                                                                                    & \multicolumn{1}{l|}{249--264}                & 248--265                                   \\ \hline
	\end{tabular}
	\caption{Periodos de siembra}
\end{table}

% Please add the following required packages to your document preamble:
% \usepackage{multirow}
\begin{table}[]
	\centering
	\begin{tabular}{llllll}
		\hline
		\multicolumn{6}{l}{\textbf{Altitude}}                                                                                                                                                                                                                                                                                                                        \\ \hline
		\multicolumn{1}{l|}{\multirow{3}{*}{\textbf{$\nu$}}} & \multicolumn{1}{l|}{\multirow{3}{*}{\textbf{\begin{tabular}[c]{@{}l@{}}Sowing window\\ $115.6 \pm 7.12\nu$\end{tabular}}}} & \multicolumn{4}{c}{\textbf{Jaccard index}}                                                                                                                         \\ \cline{3-6} 
		\multicolumn{1}{l|}{}                                   & \multicolumn{1}{l|}{}                                                                                                         & \multicolumn{2}{c|}{\textbf{Euclidean}}                                                    & \multicolumn{2}{c}{\textbf{Manhattan}}                                \\ \cline{3-6} 
		\multicolumn{1}{l|}{}                                   & \multicolumn{1}{l|}{}                                                                                                         & \multicolumn{1}{l|}{\textbf{Raw prediction}} & \multicolumn{1}{l|}{\textbf{CI prediction}} & \multicolumn{1}{l|}{\textbf{Raw prediction}} & \textbf{CI prediction} \\ \hline
		\multicolumn{1}{l|}{1}                                  & \multicolumn{1}{l|}{108 -- 122}                                                                                               & \multicolumn{1}{l|}{0.0131}                  & \multicolumn{1}{l|}{0.0568}                 & \multicolumn{1}{l|}{0.0000}                  & 0.000                  \\ \hline
		\multicolumn{1}{l|}{2}                                  & \multicolumn{1}{l|}{101 -- 129}                                                                                               & \multicolumn{1}{l|}{0.1375}                  & \multicolumn{1}{l|}{0.1888}                 & \multicolumn{1}{l|}{0.0983}                  & 0.1076                 \\ \hline
		\multicolumn{1}{l|}{3}                                  & \multicolumn{1}{l|}{94 -- 136}                                                                                                & \multicolumn{1}{l|}{0.2068}                  & \multicolumn{1}{l|}{0.3010}                 & \multicolumn{1}{l|}{0.1911}                  & 0.1944                 \\ \hline
		&                                                                                                                               &                                              &                                             &                                              &                        \\ \hline
		\multicolumn{6}{l}{\textbf{Latitude}}                                                                                                                                                                                                                                                                                                                        \\ \hline
		\multicolumn{1}{l|}{\multirow{3}{*}{\textbf{$\nu$}}} & \multicolumn{1}{l|}{\multirow{3}{*}{\textbf{\begin{tabular}[c]{@{}l@{}}Sowing window\\ $116.2 \pm 7.08\nu$\end{tabular}}}} & \multicolumn{4}{c}{\textbf{Jaccard index}}                                                                                                                         \\ \cline{3-6} 
		\multicolumn{1}{l|}{}                                   & \multicolumn{1}{l|}{}                                                                                                         & \multicolumn{2}{l|}{\textbf{Euclidean}}                                                    & \multicolumn{2}{c}{\textbf{Manhattan}}                                \\ \cline{3-6} 
		\multicolumn{1}{l|}{}                                   & \multicolumn{1}{l|}{}                                                                                                         & \multicolumn{1}{l|}{\textbf{Raw prediction}} & \multicolumn{1}{l|}{\textbf{CI prediction}} & \multicolumn{1}{l|}{\textbf{Raw prediction}} & \textbf{CI prediction} \\ \hline
		\multicolumn{1}{l|}{1}                                  & \multicolumn{1}{l|}{109 -- 123}                                                                                               & \multicolumn{1}{l|}{0.0131}                  & \multicolumn{1}{l|}{0.0568}                 & \multicolumn{1}{l|}{0.0000}                  & 0.0000                 \\ \hline
		\multicolumn{1}{l|}{2}                                  & \multicolumn{1}{l|}{102 -- 130}                                                                                               & \multicolumn{1}{l|}{0.1234}                  & \multicolumn{1}{l|}{0.1888}                 & \multicolumn{1}{l|}{0.0806}                  & 0.0909                 \\ \hline
		\multicolumn{1}{l|}{3}                                  & \multicolumn{1}{l|}{94 -- 137}                                                                                                & \multicolumn{1}{l|}{0.2183}                  & \multicolumn{1}{l|}{0.3118}                 & \multicolumn{1}{l|}{0.1884}                  & 0.1917                 \\ \hline
		&                                                                                                                               &                                              &                                             &                                              &                        \\ \hline
		\multicolumn{6}{l}{\textbf{Longitude}}                                                                                                                                                                                                                                                                                                                       \\ \hline
		\multicolumn{1}{l|}{\multirow{3}{*}{\textbf{$\nu$}}} & \multicolumn{1}{l|}{\multirow{3}{*}{\textbf{\begin{tabular}[c]{@{}l@{}}Sowing window\\ $116.2 \pm 7.3\nu$\end{tabular}}}}  & \multicolumn{4}{c}{\textbf{Jaccard index}}                                                                                                                         \\ \cline{3-6} 
		\multicolumn{1}{l|}{}                                   & \multicolumn{1}{l|}{}                                                                                                         & \multicolumn{2}{l|}{\textbf{Euclidean}}                                                    & \multicolumn{2}{c}{\textbf{Manhattan}}                                \\ \cline{3-6} 
		\multicolumn{1}{l|}{}                                   & \multicolumn{1}{l|}{}                                                                                                         & \multicolumn{1}{l|}{\textbf{Raw prediction}} & \multicolumn{1}{l|}{\textbf{CI prediction}} & \multicolumn{1}{l|}{\textbf{Raw prediction}} & \textbf{CI prediction} \\ \hline
		\multicolumn{1}{l|}{1}                                  & \multicolumn{1}{l|}{108 -- 123}                                                                                               & \multicolumn{1}{l|}{0.0263}                  & \multicolumn{1}{l|}{0.0681}                 & \multicolumn{1}{l|}{0.0000}                  & 0.0000                 \\ \hline
		\multicolumn{1}{l|}{2}                                  & \multicolumn{1}{l|}{101 -- 130}                                                                                               & \multicolumn{1}{l|}{0.1358}                  & \multicolumn{1}{l|}{0.2000}                 & \multicolumn{1}{l|}{0.0967}                  & 0.1060                 \\ \hline
		\multicolumn{1}{l|}{3}                                  & \multicolumn{1}{l|}{94 -- 138}                                                                                                & \multicolumn{1}{l|}{0.2159}                  & \multicolumn{1}{l|}{0.3225}                 & \multicolumn{1}{l|}{0.1857}                  & 0.1891                 \\ \hline
	\end{tabular}
	\caption{Comparaciones para Brandemburgo.}
\end{table}




% Please add the following required packages to your document preamble:
% \usepackage{multirow}
\begin{table}[]
	\centering
	\begin{tabular}{llllll}
		\hline
		\multicolumn{6}{l}{\textbf{Altitude}}                                                                                                                                                                                                                                                                                                                        \\ \hline
		\multicolumn{1}{l|}{\multirow{3}{*}{\textbf{$\nu$}}} & \multicolumn{1}{l|}{\multirow{3}{*}{\textbf{\begin{tabular}[c]{@{}l@{}}Sowing window\\ $115.6 \pm 7.12\nu$\end{tabular}}}} & \multicolumn{4}{c}{\textbf{Jaccard index}}                                                                                                                         \\ \cline{3-6} 
		\multicolumn{1}{l|}{}                                   & \multicolumn{1}{l|}{}                                                                                                         & \multicolumn{2}{c|}{\textbf{Euclidean}}                                                    & \multicolumn{2}{c}{\textbf{Manhattan}}                                \\ \cline{3-6} 
		\multicolumn{1}{l|}{}                                   & \multicolumn{1}{l|}{}                                                                                                         & \multicolumn{1}{l|}{\textbf{Raw prediction}} & \multicolumn{1}{l|}{\textbf{CI prediction}} & \multicolumn{1}{l|}{\textbf{Raw prediction}} & \textbf{CI prediction} \\ \hline
		\multicolumn{1}{l|}{1}                                  & \multicolumn{1}{l|}{108 -- 122}                                                                                               & \multicolumn{1}{l|}{0.3260}                  & \multicolumn{1}{l|}{0.2678}                 & \multicolumn{1}{l|}{0.3061}                  & 0.2884                 \\ \hline
		\multicolumn{1}{l|}{2}                                  & \multicolumn{1}{l|}{101 -- 129}                                                                                               & \multicolumn{1}{l|}{0.5957}                  & \multicolumn{1}{l|}{0.5178}                 & \multicolumn{1}{l|}{0.5918}                  & 0.5576                 \\ \hline
		\multicolumn{1}{l|}{3}                                  & \multicolumn{1}{l|}{94 -- 136}                                                                                                & \multicolumn{1}{l|}{0.7115}                  & \multicolumn{1}{l|}{0.7678}                 & \multicolumn{1}{l|}{0.8775}                  & 0.8269                 \\ \hline
		&                                                                                                                               &                                              &                                             &                                              &                        \\ \hline
		\multicolumn{6}{l}{\textbf{Latitude}}                                                                                                                                                                                                                                                                                                                        \\ \hline
		\multicolumn{1}{l|}{\multirow{3}{*}{\textbf{$\nu$}}} & \multicolumn{1}{l|}{\multirow{3}{*}{\textbf{\begin{tabular}[c]{@{}l@{}}Sowing window\\ $116.2 \pm 7.08\nu$\end{tabular}}}} & \multicolumn{4}{c}{\textbf{Jaccard index}}                                                                                                                         \\ \cline{3-6} 
		\multicolumn{1}{l|}{}                                   & \multicolumn{1}{l|}{}                                                                                                         & \multicolumn{2}{l|}{\textbf{Euclidean}}                                                    & \multicolumn{2}{c}{\textbf{Manhattan}}                                \\ \cline{3-6} 
		\multicolumn{1}{l|}{}                                   & \multicolumn{1}{l|}{}                                                                                                         & \multicolumn{1}{l|}{\textbf{Raw prediction}} & \multicolumn{1}{l|}{\textbf{CI prediction}} & \multicolumn{1}{l|}{\textbf{Raw prediction}} & \textbf{CI prediction} \\ \hline
		\multicolumn{1}{l|}{1}                                  & \multicolumn{1}{l|}{109 -- 123}                                                                                               & \multicolumn{1}{l|}{0.3260}                  & \multicolumn{1}{l|}{0.2678}                 & \multicolumn{1}{l|}{0.3061}                  & 0.2884                 \\ \hline
		\multicolumn{1}{l|}{2}                                  & \multicolumn{1}{l|}{102 -- 130}                                                                                               & \multicolumn{1}{l|}{0.6304}                  & \multicolumn{1}{l|}{0.5178}                 & \multicolumn{1}{l|}{0.5918}                  & 0.5576                 \\ \hline
		\multicolumn{1}{l|}{3}                                  & \multicolumn{1}{l|}{94 -- 137}                                                                                                & \multicolumn{1}{l|}{0.7307}                  & \multicolumn{1}{l|}{0.7857}                 & \multicolumn{1}{l|}{0.8979}                  & 0.8461                 \\ \hline
		&                                                                                                                               &                                              &                                             &                                              &                        \\ \hline
		\multicolumn{6}{l}{\textbf{Longitude}}                                                                                                                                                                                                                                                                                                                       \\ \hline
		\multicolumn{1}{l|}{\multirow{3}{*}{\textbf{$\nu$}}} & \multicolumn{1}{l|}{\multirow{3}{*}{\textbf{\begin{tabular}[c]{@{}l@{}}Sowing window\\ $116.2 \pm 7.3\nu$\end{tabular}}}}  & \multicolumn{4}{c}{\textbf{Jaccard index}}                                                                                                                         \\ \cline{3-6} 
		\multicolumn{1}{l|}{}                                   & \multicolumn{1}{l|}{}                                                                                                         & \multicolumn{2}{l|}{\textbf{Euclidean}}                                                    & \multicolumn{2}{c}{\textbf{Manhattan}}                                \\ \cline{3-6} 
		\multicolumn{1}{l|}{}                                   & \multicolumn{1}{l|}{}                                                                                                         & \multicolumn{1}{l|}{\textbf{Raw prediction}} & \multicolumn{1}{l|}{\textbf{CI prediction}} & \multicolumn{1}{l|}{\textbf{Raw prediction}} & \textbf{CI prediction} \\ \hline
		\multicolumn{1}{l|}{1}                                  & \multicolumn{1}{l|}{108 -- 123}                                                                                               & \multicolumn{1}{l|}{0.3478}                  & \multicolumn{1}{l|}{0.2857}                 & \multicolumn{1}{l|}{0.3265}                  & 0.3076                 \\ \hline
		\multicolumn{1}{l|}{2}                                  & \multicolumn{1}{l|}{101 -- 130}                                                                                               & \multicolumn{1}{l|}{0.6170}                  & \multicolumn{1}{l|}{0.5357}                 & \multicolumn{1}{l|}{0.6122}                  & 0.5769                 \\ \hline
		\multicolumn{1}{l|}{3}                                  & \multicolumn{1}{l|}{94 -- 138}                                                                                                & \multicolumn{1}{l|}{0.7500}                  & \multicolumn{1}{l|}{0.8035}                 & \multicolumn{1}{l|}{0.8800}                  & 0.8653                 \\ \hline
	\end{tabular}
	\caption{Comparaciones para el norte de Sajonia.}
\end{table}


% Please add the following required packages to your document preamble:
% \usepackage{multirow}
\begin{table}[]
	\centering
	\begin{tabular}{llllll}
		\hline
		\multicolumn{6}{l}{\textbf{Altitude}}                                                                                                                                                                                                                                                                                                                        \\ \hline
		\multicolumn{1}{l|}{\multirow{3}{*}{\textbf{$\nu$}}} & \multicolumn{1}{l|}{\multirow{3}{*}{\textbf{\begin{tabular}[c]{@{}l@{}}Sowing window\\ $117.4 \pm 8.29\nu$\end{tabular}}}} & \multicolumn{4}{c}{\textbf{Jaccard index}}                                                                                                                         \\ \cline{3-6} 
		\multicolumn{1}{l|}{}                                   & \multicolumn{1}{l|}{}                                                                                                         & \multicolumn{2}{c|}{\textbf{Euclidean}}                                                    & \multicolumn{2}{c}{\textbf{Manhattan}}                                \\ \cline{3-6} 
		\multicolumn{1}{l|}{}                                   & \multicolumn{1}{l|}{}                                                                                                         & \multicolumn{1}{l|}{\textbf{Raw prediction}} & \multicolumn{1}{l|}{\textbf{CI prediction}} & \multicolumn{1}{l|}{\textbf{Raw prediction}} & \textbf{CI prediction} \\ \hline
		\multicolumn{1}{l|}{1}                                  & \multicolumn{1}{l|}{109 -- 125}                                                                                               & \multicolumn{1}{l|}{0.3695}                  & \multicolumn{1}{l|}{0.3035}                 & \multicolumn{1}{l|}{0.3469}                  & 0.3269                 \\ \hline
		\multicolumn{1}{l|}{2}                                  & \multicolumn{1}{l|}{100 -- 133}                                                                                               & \multicolumn{1}{l|}{0.6666}                  & \multicolumn{1}{l|}{0.6071}                 & \multicolumn{1}{l|}{0.6938}                  & 0.6538                 \\ \hline
		\multicolumn{1}{l|}{3}                                  & \multicolumn{1}{l|}{92 -- 142}                                                                                                & \multicolumn{1}{l|}{0.8301}                  & \multicolumn{1}{l|}{0.9107}                 & \multicolumn{1}{l|}{0.8867}                  & 0.9433                 \\ \hline
		&                                                                                                                               &                                              &                                             &                                              &                        \\ \hline
		\multicolumn{6}{l}{\textbf{Latitude}}                                                                                                                                                                                                                                                                                                                        \\ \hline
		\multicolumn{1}{l|}{\multirow{3}{*}{\textbf{$\nu$}}} & \multicolumn{1}{l|}{\multirow{3}{*}{\textbf{\begin{tabular}[c]{@{}l@{}}Sowing window\\ $116.7 \pm 8.23\nu$\end{tabular}}}} & \multicolumn{4}{c}{\textbf{Jaccard index}}                                                                                                                         \\ \cline{3-6} 
		\multicolumn{1}{l|}{}                                   & \multicolumn{1}{l|}{}                                                                                                         & \multicolumn{2}{l|}{\textbf{Euclidean}}                                                    & \multicolumn{2}{c}{\textbf{Manhattan}}                                \\ \cline{3-6} 
		\multicolumn{1}{l|}{}                                   & \multicolumn{1}{l|}{}                                                                                                         & \multicolumn{1}{l|}{\textbf{Raw prediction}} & \multicolumn{1}{l|}{\textbf{CI prediction}} & \multicolumn{1}{l|}{\textbf{Raw prediction}} & \textbf{CI prediction} \\ \hline
		\multicolumn{1}{l|}{1}                                  & \multicolumn{1}{l|}{108 -- 124}                                                                                               & \multicolumn{1}{l|}{0.3695}                  & \multicolumn{1}{l|}{0.3035}                 & \multicolumn{1}{l|}{0.3469}                  & 0.3269                 \\ \hline
		\multicolumn{1}{l|}{2}                                  & \multicolumn{1}{l|}{100 -- 133}                                                                                               & \multicolumn{1}{l|}{0.6666}                  & \multicolumn{1}{l|}{0.6071}                 & \multicolumn{1}{l|}{0.6938}                  & 0.6538                 \\ \hline
		\multicolumn{1}{l|}{3}                                  & \multicolumn{1}{l|}{92 -- 141}                                                                                                & \multicolumn{1}{l|}{0.8113}                  & \multicolumn{1}{l|}{0.8928}                 & \multicolumn{1}{l|}{0.9038}                  & 0.9245                 \\ \hline
		&                                                                                                                               &                                              &                                             &                                              &                        \\ \hline
		\multicolumn{6}{l}{\textbf{Longitude}}                                                                                                                                                                                                                                                                                                                       \\ \hline
		\multicolumn{1}{l|}{\multirow{3}{*}{\textbf{$\nu$}}} & \multicolumn{1}{l|}{\multirow{3}{*}{\textbf{\begin{tabular}[c]{@{}l@{}}Sowing window\\ $116.2 \pm 7.3\nu$\end{tabular}}}}  & \multicolumn{4}{c}{\textbf{Jaccard index}}                                                                                                                         \\ \cline{3-6} 
		\multicolumn{1}{l|}{}                                   & \multicolumn{1}{l|}{}                                                                                                         & \multicolumn{2}{l|}{\textbf{Euclidean}}                                                    & \multicolumn{2}{c}{\textbf{Manhattan}}                                \\ \cline{3-6} 
		\multicolumn{1}{l|}{}                                   & \multicolumn{1}{l|}{}                                                                                                         & \multicolumn{1}{l|}{\textbf{Raw prediction}} & \multicolumn{1}{l|}{\textbf{CI prediction}} & \multicolumn{1}{l|}{\textbf{Raw prediction}} & \textbf{CI prediction} \\ \hline
		\multicolumn{1}{l|}{1}                                  & \multicolumn{1}{l|}{108 -- 123}                                                                                               & \multicolumn{1}{l|}{0.3478}                  & \multicolumn{1}{l|}{0.2857}                 & \multicolumn{1}{l|}{0.3265}                  & 0.3076                 \\ \hline
		\multicolumn{1}{l|}{2}                                  & \multicolumn{1}{l|}{101 -- 130}                                                                                               & \multicolumn{1}{l|}{0.6170}                  & \multicolumn{1}{l|}{0.5357}                 & \multicolumn{1}{l|}{0.6122}                  & 0.5769                 \\ \hline
		\multicolumn{1}{l|}{3}                                  & \multicolumn{1}{l|}{94 -- 138}                                                                                                & \multicolumn{1}{l|}{0.7500}                  & \multicolumn{1}{l|}{0.8035}                 & \multicolumn{1}{l|}{0.8800}                  & 0.8653                 \\ \hline
	\end{tabular}
	\caption{Comparaciones para el sur de Sajonia.}
\end{table}


% Please add the following required packages to your document preamble:
% \usepackage{multirow}
\begin{table}[]
	\centering
	\begin{tabular}{llllll}
		\hline
		\multicolumn{6}{l}{\textbf{Altitude}}                                                                                                                                                                                                                                                                                                                        \\ \hline
		\multicolumn{1}{l|}{\multirow{3}{*}{\textbf{$\nu$}}} & \multicolumn{1}{l|}{\multirow{3}{*}{\textbf{\begin{tabular}[c]{@{}l@{}}Sowing window\\ $117.4 \pm 8.29\nu$\end{tabular}}}} & \multicolumn{4}{c}{\textbf{Jaccard index}}                                                                                                                         \\ \cline{3-6} 
		\multicolumn{1}{l|}{}                                   & \multicolumn{1}{l|}{}                                                                                                         & \multicolumn{2}{c|}{\textbf{Euclidean}}                                                    & \multicolumn{2}{c}{\textbf{Manhattan}}                                \\ \cline{3-6} 
		\multicolumn{1}{l|}{}                                   & \multicolumn{1}{l|}{}                                                                                                         & \multicolumn{1}{l|}{\textbf{Raw prediction}} & \multicolumn{1}{l|}{\textbf{CI prediction}} & \multicolumn{1}{l|}{\textbf{Raw prediction}} & \textbf{CI prediction} \\ \hline
		\multicolumn{1}{l|}{1}                                  & \multicolumn{1}{l|}{109 -- 125}                                                                                               & \multicolumn{1}{l|}{0.4074}                  & \multicolumn{1}{l|}{0.3666}                 & \multicolumn{1}{l|}{0.3666}                  & 0.4117                 \\ \hline
		\multicolumn{1}{l|}{2}                                  & \multicolumn{1}{l|}{100 -- 133}                                                                                               & \multicolumn{1}{l|}{0.5714}                  & \multicolumn{1}{l|}{0.5263}                 & \multicolumn{1}{l|}{0.5263}                  & 0.5476                 \\ \hline
		\multicolumn{1}{l|}{3}                                  & \multicolumn{1}{l|}{92 -- 142}                                                                                                & \multicolumn{1}{l|}{0.4117}                  & \multicolumn{1}{l|}{0.4705}                 & \multicolumn{1}{l|}{0.4705}                  & 0.5769                 \\ \hline
		&                                                                                                                               &                                              &                                             &                                              &                        \\ \hline
		\multicolumn{6}{l}{\textbf{Latitude}}                                                                                                                                                                                                                                                                                                                        \\ \hline
		\multicolumn{1}{l|}{\multirow{3}{*}{\textbf{$\nu$}}} & \multicolumn{1}{l|}{\multirow{3}{*}{\textbf{\begin{tabular}[c]{@{}l@{}}Sowing window\\ $116.7 \pm 8.23\nu$\end{tabular}}}} & \multicolumn{4}{c}{\textbf{Jaccard index}}                                                                                                                         \\ \cline{3-6} 
		\multicolumn{1}{l|}{}                                   & \multicolumn{1}{l|}{}                                                                                                         & \multicolumn{2}{l|}{\textbf{Euclidean}}                                                    & \multicolumn{2}{c}{\textbf{Manhattan}}                                \\ \cline{3-6} 
		\multicolumn{1}{l|}{}                                   & \multicolumn{1}{l|}{}                                                                                                         & \multicolumn{1}{l|}{\textbf{Raw prediction}} & \multicolumn{1}{l|}{\textbf{CI prediction}} & \multicolumn{1}{l|}{\textbf{Raw prediction}} & \textbf{CI prediction} \\ \hline
		\multicolumn{1}{l|}{1}                                  & \multicolumn{1}{l|}{108 -- 124}                                                                                               & \multicolumn{1}{l|}{0.4615}                  & \multicolumn{1}{l|}{0.4137}                 & \multicolumn{1}{l|}{0.4137}                  & 0.4545                 \\ \hline
		\multicolumn{1}{l|}{2}                                  & \multicolumn{1}{l|}{100 -- 133}                                                                                               & \multicolumn{1}{l|}{0.5714}                  & \multicolumn{1}{l|}{0.5263}                 & \multicolumn{1}{l|}{0.5263}                  & 0.5476                 \\ \hline
		\multicolumn{1}{l|}{3}                                  & \multicolumn{1}{l|}{92 -- 141}                                                                                                & \multicolumn{1}{l|}{0.4200}                  & \multicolumn{1}{l|}{0.4800}                 & \multicolumn{1}{l|}{0.4800}                  & 0.5882                 \\ \hline
		&                                                                                                                               &                                              &                                             &                                              &                        \\ \hline
		\multicolumn{6}{l}{\textbf{Longitude}}                                                                                                                                                                                                                                                                                                                       \\ \hline
		\multicolumn{1}{l|}{\multirow{3}{*}{\textbf{$\nu$}}} & \multicolumn{1}{l|}{\multirow{3}{*}{\textbf{\begin{tabular}[c]{@{}l@{}}Sowing window\\ $116.2 \pm 7.3\nu$\end{tabular}}}}  & \multicolumn{4}{c}{\textbf{Jaccard index}}                                                                                                                         \\ \cline{3-6} 
		\multicolumn{1}{l|}{}                                   & \multicolumn{1}{l|}{}                                                                                                         & \multicolumn{2}{l|}{\textbf{Euclidean}}                                                    & \multicolumn{2}{c}{\textbf{Manhattan}}                                \\ \cline{3-6} 
		\multicolumn{1}{l|}{}                                   & \multicolumn{1}{l|}{}                                                                                                         & \multicolumn{1}{l|}{\textbf{Raw prediction}} & \multicolumn{1}{l|}{\textbf{CI prediction}} & \multicolumn{1}{l|}{\textbf{Raw prediction}} & \textbf{CI prediction} \\ \hline
		\multicolumn{1}{l|}{1}                                  & \multicolumn{1}{l|}{108 -- 123}                                                                                               & \multicolumn{1}{l|}{0.4800}                  & \multicolumn{1}{l|}{0.4285}                 & \multicolumn{1}{l|}{0.4285}                  & 0.4687                 \\ \hline
		\multicolumn{1}{l|}{2}                                  & \multicolumn{1}{l|}{101 -- 130}                                                                                               & \multicolumn{1}{l|}{0.5937}                  & \multicolumn{1}{l|}{0.5428}                 & \multicolumn{1}{l|}{0.5428}                  & 0.5641                 \\ \hline
		\multicolumn{1}{l|}{3}                                  & \multicolumn{1}{l|}{94 -- 138}                                                                                                & \multicolumn{1}{l|}{0.4666}                  & \multicolumn{1}{l|}{0.5333}                 & \multicolumn{1}{l|}{0.5333}                  & 0.6170                 \\ \hline
	\end{tabular}
	\caption{Comparaciones para Baviera.}
\end{table}


\newpage
\newpage
\section{Cronograma de actividades}



Durante el semestre en curso se contempla emplear los datos correspondientes a las fechas de siembra de maíz en Alemania, reportados por \citet[Tabla 1]{Parker_2016}, a partir de los criterios topográficos y geográficos descritos por \cite{Flores_2025}, ver Tabla \ref{Real_data}. Los cuales, tras identificar una métrica apropiada, sirvan para realizar una comparación complementaria con resultados previamente publicados.\\

\begin{table}[h]
	\centering
	\begin{tabular}{l|c|c}
		\hline
		\multicolumn{1}{c|}{\textbf{Región}} & \textbf{Características} & \textbf{Fecha de siembra} \\ \hline
		
		\multirow{3}{*}{\parbox{4.5cm}{\textbf{Brandemburgo}\\ \textbf{Norte de Sajonia}}} 
		& Altitud menor que $250$ m.s.n.m. 
		& $115.6 \pm 7.12$ \\ \cline{2-3}
		
		
		& Latitud mayor o igual que $51^\circ$ N. 
		& $116.2 \pm 7.08$ \\ \cline{2-3}
		
		& Longitud mayor o igual que $8.5^\circ$ E. 
		& $116.2 \pm 7.3$ \\ \hline
		
		
		\multirow{3}{*}{\parbox{4.5cm}{\textbf{Sur de Sajonia}\\ \textbf{Baviera}}}  
		& Altitud mayor o igual que $250$ m.s.n.m. 
		& $117.4 \pm 8.29$ \\ \cline{2-3}
		
		
		& Latitud menor que $51^\circ$ N. 
		& $116.7 \pm 8.23$ \\ \cline{2-3}
		
		& Longitud mayor o igual que $8.5^\circ$ E. 
		& $116.2 \pm 7.3$ \\ \hline
		
	\end{tabular}
	\caption{Fechas promedio de siembra de maíz y sus correspondientes desviaciones estándar reportadas por \cite{Parker_2016}, asociadas a las características geográficas y topográficas descritas por \cite{Flores_2025} para Brandemburgo, Sajonia y Baviera.}
	\label{Real_data}
\end{table}


Adicionalmente, se llevará a cabo la construcción de los intervalos de confianza correspondientes mediante los métodos de bootstrap y jackknife, así como la integración de restantes en la redacción del trabajo en formato de artículo. Finalmente, se contempla la presentación de los avances y resultados en el Congreso Internacional de Ecología Estadística (ISEC, por sus siglas en inglés), que se llevará a cabo del 10 al 15 de enero de 2027 en Mérida, México,   \href{https://drive.google.com/file/d/1y_0bCZc6UKrdhD0rQ020FxtoEe3sbJ3-/view?usp=sharing}{ver comprobante de aceptación aquí}.



\begin{figure}[h]
	\resizebox{\linewidth}{!}{%
		\centering
		\begin{ganttchart}[
			hgrid,
			vgrid,
			x unit=1.8cm,
			y unit chart=0.7cm,
			bar height=0.5,
			group height=0.7,
			title height=1,
			milestone height=.6,
			bar/.style={fill=blue!50},
			group/.style={fill=gray!50},
			title/.style={fill=gray!20},
			milestone/.style={fill=red},
			bar label font=\small,
			group label font=\bfseries\small
			]{1}{6}
			
			% -------------------------------------------------
			% ENCABEZADOS
			% -------------------------------------------------
			
			\gantttitle{Cronograma de actividades}{6} \\
			
			\gantttitle{Ago. 2026}{1}
			\gantttitle{Sept. 2026}{1}
			\gantttitle{Oct. 2026}{1}
			\gantttitle{Nov. 2026}{1}
			\gantttitle{Dic. 2026}{1}
			\gantttitle{Ene. 2027}{1}
			\\
			
			% -------------------------------------------------
			% ANÁLISIS DE DATOS
			% -------------------------------------------------
			
			\ganttgroup{Análisis de datos y comparación}{1}{5} \\

			\ganttbar
			{Identificación de una métrica apropiada}
			{2}{3} \\
			
			\ganttbar
			{Comparación con resultados previamente publicados}
			{3}{5} \\
			
			% -------------------------------------------------
			% INTERVALOS DE CONFIANZA
			% -------------------------------------------------
			
			\ganttgroup{Construcción de intervalos de confianza}{1}{2} \\
			
			\ganttbar
			{Bootstrap}
			{1}{2} \\
			
			\ganttbar
			{Jackknife}
			{1}{2} \\
			
			% -------------------------------------------------
			% PRODUCCIÓN CIENTÍFICA
			% -------------------------------------------------
			
			\ganttgroup{Integración y redacción científica}{4}{5} \\
			
			\ganttbar
			{Integración de resultados en el artículo}
			{4}{5} \\
			
			% -------------------------------------------------
			% PRESENTACIÓN
			% -------------------------------------------------
			
			\ganttgroup{Presentación de resultados}{5}{6} \\
			
			\ganttbar
			{Preparación de resultados para el ISEC}
			{5}{5} \\
			
			\ganttmilestone
			{ISEC 2027}
			{6}
			
		\end{ganttchart}
		
	} % fin resizebox
	
	\caption{Cronograma de actividades correspondientes al semestre agosto--diciembre de 2026.}
	\label{fig:gantt_semestre}
	
\end{figure}